\documentclass[11pt]{article}
\usepackage{fontspec,amsmath,amssymb,amsthm,hyperref}
\usepackage[margin=1in]{geometry}
\setlength{\emergencystretch}{3em}
\title{Compatibility exports for the independent product realization}
\author{LangevinLean mathematical exposition}
\date{}
\begin{document}\maketitle
\noindent Source module: \texttt{IPLA/ProductReference.lean}. Compile this proof document with XeLaTeX or LuaLaTeX. The accompanying HTML pairs the mathematical argument with the exact Lean statements and records the scope of the formalization.


This file re-exports the shared product realization and its law-identification theorem. A product of an initial coupling and a noise source realizes every needed coupling while fixing a common algorithm output law.


\section{Separate initialization from the noise source}


On \((H\times H)\times S\) with law \(c\otimes Q\), let the two initial variables be the first and second coordinates of \(c\), and let every innovation depend only on \(S\). Product independence makes both initial variables independent of the noise source. The first chain evolves by \(x\mapsto x-h\nabla V(x)+\sqrt{2h}\xi_k\), and its output law is the pushforward of \(\mu_0\otimes Q\) by the measurable recursion; hence it depends only on \(\mu_0\), not on the chosen coupling \(c\). This verifies the common-output obligation needed to take the coupling infimum. Finite coupling cost and Gibbs moments supply the \(L^2\) initial realization; the exported output-finiteness statement guards subsequent real transport bounds.



The \texttt{export} commands expose the already proved declarations under the IPLA namespace. Their definitions and proofs are unchanged, so the mathematical argument is the one in the imported source modules; the dependency diagram below opens those full proofs. This file itself makes no additional analytic assumption and contains no replacement proof.


\end{document}
